\documentclass[11pt,a4paper]{article}
\usepackage[margin=2.5cm]{geometry}
\usepackage{xcolor}
\usepackage{titlesec}
\usepackage{enumitem}
\usepackage{fancyhdr}
\usepackage{hyperref}
\usepackage{parskip}

% PJD Brand Colours
\definecolor{pjdnavy}{HTML}{0A192F}
\definecolor{pjdgold}{HTML}{FFD700}

% Section formatting
\titleformat{\section}{\large\bfseries\color{pjdnavy}}{}{0em}{}[\color{pjdgold}\titlerule]
\titleformat{\subsection}{\normalsize\bfseries\color{pjdnavy}}{}{0em}{}

% Header/Footer
\pagestyle{fancy}
\fancyhf{}
\fancyhead[L]{\small\color{pjdnavy}\textbf{Paul Jeggels Designs}}
\fancyhead[R]{\small\color{pjdnavy}Surfboard Care \& Maintenance FAQ}
\fancyfoot[C]{\small\color{pjdnavy}15 Dageraad Street, Jeffreys Bay \,|\, +27 82 960 9353 \,|\, pauljeggelsdesigns.co.za}
\fancyfoot[R]{\small\thepage}
\renewcommand{\headrulewidth}{0.4pt}
\renewcommand{\footrulewidth}{0.4pt}

% Q&A formatting
\newcommand{\question}[1]{\subsection*{Q: #1}}
\newcommand{\answer}{\noindent\textbf{A:} }

\begin{document}

\begin{center}
{\LARGE\bfseries\color{pjdnavy} Paul Jeggels Designs}\\[6pt]
{\Large\color{pjdnavy} Surfboard Care \& Maintenance FAQ}\\[4pt]
{\small\color{gray} Keep Your Board in Top Shape --- Tips from 40+ Years of Shaping}
\end{center}

\vspace{1em}

% ─────────────────────────────────────────────────────────────────────
\section{Daily Care}
% ─────────────────────────────────────────────────────────────────────

\question{What should I do after every surf session?}
\answer Three simple things:
\begin{enumerate}[leftmargin=1.5em]
    \item \textbf{Rinse with fresh water.} Salt water corrodes fin screws, breaks down wax unevenly, and dries out the glass over time. A quick hose-down after each session goes a long way.
    \item \textbf{Dry it off.} Wipe down the board and let it air-dry in the shade before putting it in a bag or on a rack. Never seal a wet board in a board bag --- trapped moisture leads to delamination.
    \item \textbf{Store it out of direct sunlight.} UV radiation is the number one killer of surfboards. Even 30 minutes in the Eastern Cape sun can yellow the resin and weaken the glass.
\end{enumerate}

\question{Does salt water damage my board?}
\answer Not immediately, but over time, yes. Salt crystals can work their way into micro-cracks and tiny dings you haven't noticed. If salt water gets under the glass (through a ding), it saturates the foam core. That's when you get waterlogged boards that feel heavy and sluggish. The fix is simple: rinse after every surf and repair dings quickly.

% ─────────────────────────────────────────────────────────────────────
\section{Wax Maintenance}
% ─────────────────────────────────────────────────────────────────────

\question{How often should I wax my board?}
\answer Apply a fresh top coat every 2--3 sessions or whenever the wax feels smooth and slippery. Every month or so (more often in summer), strip all the old wax off completely and start fresh. Old wax builds up, goes hard, and stops gripping properly. You'll notice the difference immediately after a full re-wax.

\question{How do I strip old wax?}
\answer Leave your board in the sun for 5--10 minutes (no longer!) to soften the wax. Then use a wax comb to scrape it off. For the last residue, a bit of citrus-based wax remover or even flour rubbed on the surface works well. Wipe clean with a cloth. Don't use petrol or harsh solvents --- they can damage the resin.

\question{What type of wax should I use in Jeffreys Bay?}
\answer J-Bay water temperature ranges from about 16\textdegree C in winter to 22\textdegree C in summer. Use \textbf{cool water wax} for most of the year. In the warmer months (December--March), you can switch to \textbf{warm water wax} for a harder base that won't go mushy in the heat. Apply a base coat first with a harder wax, then a top coat with your temperature-appropriate wax.

% ─────────────────────────────────────────────────────────────────────
\section{Dings and Damage}
% ─────────────────────────────────────────────────────────────────────

\question{How do I check for dings?}
\answer Run your hand slowly over the entire board after each session, especially the rails (edges), nose, and tail. Look for:
\begin{itemize}[leftmargin=1.5em]
    \item \textbf{Cracks}: Visible lines in the glass, often near the fin boxes or around the nose
    \item \textbf{Pressure dents}: Soft spots where the foam has compressed (usually where your feet go)
    \item \textbf{Chips}: Small pieces of fibreglass broken away, exposing the foam underneath
    \item \textbf{Delamination}: Areas where the glass has separated from the foam --- sounds hollow when you tap it
\end{itemize}
Any crack or chip that exposes the white foam underneath needs to be sealed before you surf again. Water getting into the foam is the real problem.

\question{Can I fix small dings myself?}
\answer Minor chips and small cracks (under 2cm) can be temporarily sealed with Solarez UV resin --- a tube you can buy at most surf shops. Squeeze it on, smooth it out, and leave it in the sun for 3--5 minutes to cure. It's not a permanent fix, but it stops water getting in until you can bring it to Paul for a proper repair.

For anything bigger --- rail damage, cracked fins, nose breaks, or delamination --- bring it to the workshop. A bad DIY repair can make the problem worse and cost more to fix later.

\question{What happens if water gets into my board?}
\answer Water-logged foam adds weight and kills performance. Your board will feel sluggish, heavy to carry, and slow in the water. The foam can also start to rot if it stays wet for a long time, especially in warmer weather. If you suspect water damage (the board feels heavier than usual or you can see discolouration through the glass), bring it to Paul. He'll assess whether it needs to be drained, dried, and re-glassed or whether the damage is too far gone.

\question{How much does a ding repair cost at PJD?}
\answer Ding repairs start from R500 for minor work. The price depends on the size and location of the damage. A small chip or crack is quick and affordable. A full rail rebuild or nose repair takes more time and materials. Paul and the team will give you an honest assessment --- if it's a 5-minute fix, you won't be charged for an hour's work.

% ─────────────────────────────────────────────────────────────────────
\section{Sun, Heat \& Storage}
% ─────────────────────────────────────────────────────────────────────

\question{Can I leave my board in the car?}
\answer Not for long. A car interior in the South African sun can reach 60--70\textdegree C. At those temperatures, the foam inside your board can expand and separate from the fibreglass (delamination). You'll hear it --- a cracking sound --- and see bubbles forming under the glass. If you must leave your board in the car, crack the windows, cover it with a reflective board bag, and park in the shade. Better yet, take it out of the car.

\question{What's the best way to store my board at home?}
\answer Best options, in order:
\begin{enumerate}[leftmargin=1.5em]
    \item \textbf{Wall rack indoors} --- horizontal or at a slight angle, fins up, out of direct sunlight. This is ideal.
    \item \textbf{Board bag indoors} --- good protection from dust and light knocks. Make sure the board is dry before you bag it.
    \item \textbf{Covered outdoor rack} --- fine if it's shaded and sheltered from rain and direct sun.
\end{enumerate}
Avoid storing your board standing upright on its tail for long periods --- the weight can cause pressure dents in the tail area over time. Horizontal is always better.

\question{Does UV damage affect performance?}
\answer Yes. UV radiation yellows the resin (cosmetic) and over time weakens the fibreglass (structural). A board that's been left in the sun regularly will become more brittle and prone to cracking. Paul uses UV-stable resin in his builds, which helps, but no resin is completely immune to prolonged sun exposure. The best protection is keeping your board covered or indoors when you're not using it.

% ─────────────────────────────────────────────────────────────────────
\section{Fins}
% ─────────────────────────────────────────────────────────────────────

\question{What's the difference between FCS and Futures fin systems?}
\answer Both are excellent systems --- the difference is in how the fins attach:
\begin{itemize}[leftmargin=1.5em]
    \item \textbf{FCS II}: Two-tab system. Fins click in without screws (though you can add a screw for extra security). Easy to swap fins on the beach.
    \item \textbf{Futures}: Single-tab system with one screw per fin. Some surfers feel Futures have a more solid, locked-in connection. Slightly more effort to change.
\end{itemize}
Paul installs both systems. If you don't have a preference, he'll recommend based on the board design and your surfing style. Neither is objectively better --- it's personal preference and what fins you already own.

\question{How do I look after my fins?}
\answer Rinse them with fresh water along with the board. Check fin screws regularly --- they can loosen over time, especially with FCS systems. If a fin feels wobbly, tighten the grub screw with the fin key (a small Allen key). Don't over-tighten --- just snug. If the fin box itself feels loose or cracked, bring it to Paul. Fin box repairs are straightforward but important to get right.

\question{When should I change my fins?}
\answer If a fin is chipped, cracked, or has lost its foil shape, it's affecting your surfing. Small chips on the leading edge can be sanded smooth. Larger damage means it's time for a replacement. If you want to experiment with different fin templates (more drive, more release, more hold), Paul manufactures custom fins matched to your board. Stock fins work fine, but custom fins matched to your specific board's flex and rocker can make a noticeable difference.

% ─────────────────────────────────────────────────────────────────────
\section{When to See Paul}
% ─────────────────────────────────────────────────────────────────────

\question{How do I know when a board needs professional repair vs a DIY fix?}
\answer Simple rule: if you can see white foam, and the damaged area is bigger than a 50c coin, bring it in. Also bring it in if:
\begin{itemize}[leftmargin=1.5em]
    \item The damage is on the rail (structural area)
    \item The fin box is cracked or loose
    \item The nose or tail is broken
    \item You can feel soft, spongy areas (water damage)
    \item The board has delamination (glass separating from foam)
\end{itemize}
DIY patches are fine as a temporary seal to keep water out, but they're not designed to restore structural integrity. Paul's repairs are done with proper materials and last the life of the board.

\question{Where is the PJD workshop for repairs?}
\answer 15 Dageraad Street, Jeffreys Bay. You can drop your board off or give Paul a call first on +27 82 960 9353 to discuss the damage. Most minor repairs are turned around in a few days. Bigger jobs may take a week or so depending on how busy the workshop is.

\vfill
\begin{center}
\small\color{gray}
Paul Jeggels Designs --- 15 Dageraad Street, Jeffreys Bay, South Africa\\
Phone: +27 82 960 9353 --- Hand-shaped surfboards since the 1980s
\end{center}

\end{document}
